\section{Introduction}
UAV vehicle detection must localize small targets under changing altitude, viewpoint, scale, illumination, motion, and clutter. These conditions are challenging for generic object detection and for aerial imaging in particular \cite{zhao2019object,cheng2016survey,xia2022dota,zhu2017deepRS}. RGB imagery provides texture and color, but may weaken under poor contrast or changing exposure. Thermal observations can provide temperature-related contrast, while event representations can encode brightness changes with high temporal resolution and wide dynamic range \cite{baltrusaitis2019multimodal,gallego2022event}. The key design question is therefore not merely how to concatenate modalities, but how to use their complementary information without unnecessarily increasing the cost of a lightweight detector.

A second issue is evaluation integrity. In multimodal UAV recordings, neighboring observations may be visually similar even when their file contents are not identical. If such observations fall on different sides of a split, validation performance can be optimistic. This risk is distinct from model design and should be handled before interpreting a fusion gain \cite{kapoor2023leakage,cawley2010overfitting}. In the present project, an earlier split review identified exact RGB duplicates and subsequently identified cross-partition adjacent-or-near-identical observations through a human review workflow. These findings motivate a more conservative validation protocol rather than a stronger claim about universal independence.

This paper presents RA-RepDet, a compact five-channel RGB-thermal-event detector. The reliability-aware model uses separate stems for RGB, thermal, and event inputs, predicts image-level softmax fusion weights, and passes the fused representation through a RepViT-M0.9 \cite{wang2024repvit}, feature-pyramid network (FPN) \cite{lin2017fpn}, and FCOS detector \cite{tian2019fcos}. The matched early-fusion baseline differs only in its front-end fusion mechanism. During reliability-aware training, independent modality dropout reduces reliance on a single input group; the $p=0.15$ setting was specified from archived development evidence before V40 results were viewed. V40 does not perform a dropout sweep or select an optimal dropout rate.

The contribution of the work is twofold. First, it offers a lightweight, reliability-aware fusion front end for a shared detector stack. Second, it reports the result on a validation partition constructed to keep all components of a fixed candidate-and-adjacency graph within one side of the split. The conclusions are intentionally narrow: they concern a fixed validation protocol and controlled synthetic channel removal, not independent testing or real sensor-failure deployment.

\noindent\textbf{Contributions.}
\begin{itemize}
    \item We develop a lightweight tri-modal fusion front end with modality-specific stems, image-level softmax fusion weights, and a shared RepViT-FPN-FCOS detector.
    \item We compare matched early fusion and a pre-specified reliability-aware $p=0.15$ configuration under a frozen V40-v2 expanded-adjacency component-disjoint validation protocol.
    \item We report two-seed aggregate results, image-level bootstrap intervals, deterministic synthetic channel-removal results, and a reproducible efficiency measurement with explicit latency boundaries.
    \item We document the scope and limitations of the validation evidence, including the absence of an independent test set, trained single-modality baselines, and a V40 dropout sweep.
\end{itemize}

\section{Related Work}
\subsection{Aerial Detection and Lightweight Detectors}
Aerial detection involves wide scale variation, small objects, geometric changes, and cluttered backgrounds \cite{cheng2016survey,zhu2017deepRS,xia2022dota}. Modern detectors include region-proposal and dense one-stage formulations \cite{ren2017faster,lin2020focal}. FCOS is an anchor-free dense detector that predicts object locations directly \cite{tian2019fcos}. In this study, the detector head is held fixed; the research question is whether the tri-modal fusion front end can improve a compact detection pipeline. RepViT is used as a mobile-oriented backbone, and FPN provides multi-scale features \cite{wang2024repvit,lin2017fpn}.

\subsection{Multimodal and Event-Based Perception}
Multimodal learning integrates complementary sources while accounting for differences in their reliability and representation \cite{baltrusaitis2019multimodal,ramachandram2017deep}. Visible-infrared fusion has been studied extensively at pixel and feature levels \cite{li2017pixel,ma2019surveyfusion,li2019densefuse}. Event cameras provide asynchronous brightness-change information and can support perception under fast motion or high dynamic range \cite{lichtsteiner2008sensor,gallego2022event,gehrig2021dsec}. RA-RepDet does not reconstruct raw event streams or infer a missing modality. It consumes a preconstructed five-channel sample and learns a compact image-level mixture of modality-stem features.

\subsection{Missing Modality and Evaluation Validity}
Modality dropout is a common strategy for limiting brittle dependence on one information source \cite{neverova2016moddrop}. Here, it is used as a training-time intervention, while zeroing one channel group at evaluation time gives a controlled synthetic channel-removal protocol. This protocol is useful for comparing models under the same intervention, but it cannot reproduce sensor drift, misregistration, field-of-view change, or other physical failure modes.

Evaluation leakage and adaptive reuse of validation information can produce optimistic performance estimates \cite{kapoor2023leakage,cawley2010overfitting}. The present work does not claim that any finite audit proves the absence of all dependence. Instead, it constructs a split whose train-validation boundary does not cut components of a specified graph formed from exact RGB, pHash, dHash, and human-adjudicated adjacent-or-near-identical relations.
